\section{经典习惯化特征的补充检验}
\label{sec:hallmarks}
\subsection{实验范围与总体结果}
按照Rankin等人的分类\citep{rankin2009}，在已有递减和自发恢复实验之外，进一步考察多轮训练、重复频率、输入强度、平台后追加刺激、刺激特异性和去习惯化。H1--H10分别表示这些经典特征；操作性判据见附录。

补测保持完整M1及其$U$和恢复常数不变，采用3701--3703。主实验包括三轮训练、多时点恢复、150/200Hz输入以及12次与24次训练的比较，并设置两个替代刺激B。另增加2s轮间休息条件，以及一次针对H4(b)比较条件不足的补测，设计顺序分别在下文说明。共获得525个观测窗口，其中6个为空白检查窗。

\begin{table}[!htbp]\centering\small
\caption{经典特征的实验结果。比例按输入实现计算。缺少有效比较条件时记为无法判断，未实施的实验另列。}
\label{tab:hallmarks}
\begin{tabularx}{\linewidth}{L{2.7cm}Y}\toprule
特征 & 结果\\\midrule
H1 响应递减 & 首轮M1短间隔5/5、新基准3/3符合递减判据；M0在已测条件下计数恒定。\\
H2 自发恢复 & 首轮10s探测5/5符合判据；新短间隔组5s探测3/3达到初次响应的至少80\%。\\
H3 多轮增强 & 2s部分恢复条件3/3符合补充判据；10s条件0/3符合原判据。\\
H4(a) 频率依赖递减 & 更高重复频率引起更强的相对递减；结果基于早晚期均值比较。\\
H4(b) 频率依赖恢复 & 原比较1/3满足分析条件；24次补测3/3满足，但F50均未分辨出更早区间，初次80\%仅1/3在高频下明确更早。\\
H5 强度依赖 & 150/200Hz比较方向3/3一致，2/3达到预设D差。\\
H6 平台后追加刺激 & 12次与24次训练后的五点恢复计数逐实现相同，未观察到预设效应。\\
H7 刺激特异性 & 两个B的aBN1基线均为零，未能判断。\\
H8 去习惯化 & 两个B均激活非输入神经元，但后续A响应未超过等时休息对照。\\
H9 解除效应的递减 & 无可靠H8效应，未继续测试。\\
H10 长期习惯化 & 未实施长期协议或相应生物标定。\\\bottomrule
\end{tabularx}\end{table}

新基准短间隔组的平均$D$为98.0\%，$T=1$s组为40.8\%。两者与首轮90.22\%和9.88\%的差异涉及输入实现及刺激条件变化：首轮较长间隔为2s，补测为1s，故不作跨阶段的直接效应量比较。

\subsection{多轮训练与轮间恢复程度（H3）}
原主实验每轮12次刺激，共三轮，轮间休息10s。判据要求后续各轮初次响应达到首轮至少80\%，再比较第三轮相对于首轮的递减比例或半响应到达序号。三个实现均满足初次恢复条件，但均未达到增强判据。

Smart等人的形式化分析允许在部分恢复后考察后续轮次的响应减弱\citep{smart2024}。据此，在读取主实验多轮及恢复结果之前，补充固定2s轮间休息条件；设计时已读取B实验中的单轮A基线。原10s数据及判据保留，2s条件单独分析。

补充判据要求第二、三轮的初次响应均高于前轮末期，表明轮间存在自发恢复；同时，第三轮前3次均值较首轮至少降低1脉冲，或以首轮初次为共同基准的半响应到达至少提前一次，中间轮不反向变化。三个实现均符合这一判据。

\begin{table}[!htbp]\centering\small\setlength{\tabcolsep}{3pt}
\caption{2s轮间休息下的三轮响应。半响应序号指首次连续三次不超过基准50\%的起始序号。共同基准为首轮初次响应，各轮基准为相应轮次的初次响应。}
\begin{tabular}{rllll}\toprule
种子 & 各轮首次 & 各轮前3次均值 & 共同基准半响应 & 各轮基准半响应\\\midrule
\inputrows{data/h3.tex}\bottomrule
\end{tabular}\end{table}

\newcommand{\CyclePanel}[2]{\begin{tikzpicture}\begin{axis}[width=.94\linewidth,height=4.9cm,title={#2},xlabel={本轮刺激序号},ylabel={aBN1脉冲数},xmin=1,xmax=12,ymin=0,ymax=10,xtick={1,4,8,12},grid=major,grid style={black!10},tick label style={font=\small}]
\addplot[Blue,thick,mark=*] table[x=trial,y=mean]{data/h3_#1_1.csv};
\addplot[Orange,thick,dashed,mark=triangle*] table[x=trial,y=mean]{data/h3_#1_2.csv};
\addplot[Green,thick,dashdotted,mark=square*] table[x=trial,y=mean]{data/h3_#1_3.csv};
\end{axis}\end{tikzpicture}}
\begin{figure}[!htbp]\centering
\begin{minipage}{.48\linewidth}\CyclePanel{10}{轮间10s}\end{minipage}\hfill
\begin{minipage}{.48\linewidth}\CyclePanel{2}{轮间2s}\end{minipage}
\par\smallskip\small\textcolor{Blue}{蓝：第一轮}\quad\textcolor{Orange}{橙：第二轮}\quad\textcolor{Green}{绿：第三轮}
\caption{不同轮间休息时长下的训练响应。曲线为三个实现的均值，部分曲线重合。各轮之间连续运行，未人工重置资源。}
\end{figure}

采用共同基准时，三例的半响应起点均由第3次提前至第2次；采用各轮自身初次为基准时，3703的序号为3/3/3，3702为3/2/3，未表现为一致的逐轮加快。较弱的后续响应可由轮间资源未完全恢复解释。实验中$U$和恢复常数始终固定，因此结果支持残余状态造成的多轮历史效应，不涉及学习率改变。

\subsection{重复频率与恢复速度（H4）}
更频繁刺激引起的较大递减符合H4(a)的方向。H4(b)进一步要求比较训练后的恢复速度。首轮仅设置2s和10s探测，且长间隔组递减较弱，难以据此估计恢复快慢。主补测改为$T=0.5/1$s、各12次训练，并在末态后分别等待0.2/0.5/1/2/5s探测。

预设分析以损失恢复比例$F$的50\%越阈区间，以及0.5/1/2s时的F差为主要指标。主实验开始前还规定报告响应达到初次80\%的时间区间，以考察结果对归一化方式的依赖。

\newcommand{\RecoveryPanel}[3]{\begin{tikzpicture}\begin{axis}[width=.93\linewidth,height=5cm,title={#3},xlabel={末响应窗后等待（s）},ylabel={#2},xmin=0,xmax=5.2,xtick={0,1,2,5},grid=major,grid style={black!10},tick label style={font=\small},title style={font=\small}]
\addplot[Blue,thick,mark=*] table[x=rest_s,y=#1]{data/original_fast12.csv};
\addplot[Orange,thick,dashed,mark=triangle*] table[x=rest_s,y=#1]{data/original_slow12.csv};
\end{axis}\end{tikzpicture}}
\begin{figure}[!htbp]\centering
\begin{minipage}{.48\linewidth}\RecoveryPanel{relative_initial}{$Q=R/I$}{相对初次响应}\end{minipage}\hfill
\begin{minipage}{.48\linewidth}\RecoveryPanel{F}{$F=(R-L)/(I-L)$}{损失恢复比例}\end{minipage}
\par\smallskip\small\textcolor{Blue}{蓝：$T=0.5$s}\quad\textcolor{Orange}{橙：$T=1$s}
\caption{原12次训练后恢复曲线的两种归一化表示。图中保留三个实现，线段连接采样点，未拟合恢复函数；各实现是否满足比较条件另行检验。}
\end{figure}

原比较仅3702同时满足两组的初次响应、递减和平台条件。3701较低频组末两组三次响应均值之差为0.667，大于允许值$0.10E=0.600$；3703较低频组D为29.4\%，低于30\%阈值。这两例未纳入有效恢复速度比较。

此外，$T=1$s组在0.2s或0.5s等待后探测，加上300ms响应窗，实际起始间隔为0.5s或0.8s，短于训练时的1s。因而，这些早期点反映提前到来的刺激响应，不宜单独用来解释停止重复刺激后的恢复。

看到上述结果后，仅追加一次补测：新设$T=0.75$s、24次训练，与已有$T=0.5$s、24次匹配，并在0.5/1/2/5s探测。方案在新增仿真前固定，神经参数及分析阈值不变。该设计由主结果引出，属于结果驱动的补充实验；所有探测均晚于训练节律下的下一次刺激起始。

\begin{table}[!htbp]\centering\small
\caption{24次训练后的恢复时间区间。三例均满足比较条件。区间单位为秒，表示在采样单调性假定下首次观测越阈所在范围；不表示置信区间。}
\begin{tabular}{rllll}\toprule
种子 & 高频F50 & 较低频F50 & 高频首次80\% & 较低频首次80\%\\\midrule
\inputrows{data/h4.tex}\bottomrule
\end{tabular}\end{table}

三例F50均落在相同的采样区间，未分辨出高频组更早恢复。按初次80\%评价，仅3702高频组明确较早：其较低频组5s计数为7，尚未达到初次9的80\%。因此，H4(b)未获得一致支持。区间相同仍可能包含真实时间差异，而单例越阈差异也不足以推断普遍的频率效应。

\subsection{输入强度与平台后追加刺激（H5、H6）}
H5在相同200ms刺激宽度和0.5s重复间隔下比较150Hz与200Hz人工输入率。这里改变的是单次刺激内的神经驱动率，与H4改变的重复率$1/T$不同。两组初次响应均超过5脉冲，预设较弱输入的D应高出至少0.10。

\begin{table}[!htbp]\centering\small
\caption{输入率对递减比例的影响。D以百分数表示，差值单位为百分点。人工输入率尚未对应到实际机械刺激强度。}
\begin{tabular}{rrrrl}\toprule
种子 & 150Hz D & 200Hz D & 差值 & 达到10个百分点\\\midrule
\inputrows{data/h5.tex}\bottomrule
\end{tabular}\end{table}

三个实现的方向一致，两个达到预设差值；150Hz与200Hz的平均D分别为98.0\%和83.3\%。另一例差异为8.3个百分点，未达到阈值。阶段二400ms刺激也曾出现相近方向，但当时属于事后探索；本次采用新输入实现和预先确定的比较方案，增加了对这一方向的支持。

H6考察达到输出平台后继续刺激是否延缓恢复。在第12次刺激后保存状态，再沿相同轨迹运行至第24次，分别进行五点恢复探测。三个实现的两个末态均满足平台条件，但对应探测计数差均为$[0,0,0,0,0]$，未观察到额外训练对输出恢复的影响。资源变量的微小变化尚未表现为可检测的计数差异；结论限于当前训练长度与采样分辨率。

\subsection{替代刺激的特异性与去习惯化效应（H7、H8）}
以JO-CE为原刺激A，选择JO-F（60个细胞）及上游标注\texttt{neu\_JON\_D\_m}（16个细胞）为两个候选B，均给予150Hz、200ms输入。STD仍仅施加于JO-CE输出，B刺激期间不人工重置资源。

H7要求在A响应下降后，仍能比较B是否保留正常响应。因此，预设未训练B的aBN1计数至少5脉冲，且为未训练A的0.5--2倍。两个B自身的aBN1响应均为零，未满足共同读出的基线条件，本次未能判断刺激特异性。

H8比较的是B之后原刺激A的响应是否增强，不要求B自身驱动aBN1。两个B均产生至少100个非输入神经元脉冲，B后的空白窗口aBN1为零，满足网络激活及残留输出的控制条件。随后比较训练后B$\to$A与等时休息$\to$A，并设置未训练B$\to$A，以检查一般增益。

\par\medskip\noindent\begin{minipage}{\linewidth}\centering\small\setlength{\tabcolsep}{4pt}
\captionof{table}{B刺激及后续A响应。训练增益为训练后B$\to$A相对于等时休息$\to$A的差值，未训练增益在对应初态下计算。}
\begin{tabular}{rlrrrrr}\toprule
种子 & B & 非输入脉冲 & 休息后A & B后A & 训练增益 & 未训练增益\\\midrule
\inputrows{data/b_controls.tex}\bottomrule
\end{tabular}\end{minipage}\par\medskip

所有增益均为零，表明本次所选B群体、强度与时序未诱发去习惯化。第\ref{sec:mechanism}节的资源重置虽能恢复响应，但属于直接改变内部状态，与这里的替代刺激实验不同。由于未得到可靠解除效应，H9未继续测试。H10所需的长期训练和生物时间尺度标定亦未实施。
\FloatBarrier
